% \tracingcommands NAMES WHAT A COMMAND MEANS, NOT WHAT IT IS CALLED.
%
% The line is the text \meaning would give for the token, not the token's
% own name:
%
%   \let\foo=\relax    \foo    {\relax}
%   \countdef\cc=7     \cc     {\count7}
%   \dimendef\dd=3     \dd     {\dimen3}
%   \toksdef\tk=4      \tk     {\toks4}
%   \skipdef\sk=5      \sk     {\skip5}
%   \let\bg={          \bg     {begin-group character {}
%   \let\pp=\par       \pp     {\par}
%   \chardef\ch=65     \ch     {\char"41}
%   an undefined one           {undefined}
%
% This is what trip's `\let\paR=\par' was showing: hstex drew {\paR} where
% the reference draws {\par}, three times over. It also explains the
% {undefined} and the {|relax} lines the pass-two log was missing --
% |relax being \relax with \escapechar set to `|'.
%
% ONE LINE PER COMMAND, even where the engine puts a token in the place of
% the one it read and reads it again. `\let\bg={ \bg' draws one line, not
% a {begin-group character {} for the brace on top of one for \bg.
%
% A \chardef'd control sequence, and \char, BELONG TO A RUN OF CHARACTERS
% like the characters themselves: a run is traced once, at its first, so
% `M\ch' draws {the letter M} and nothing for the \ch. Where the run has
% not started, \ch draws {\char"41} -- the meaning, not the character.
\tracingonline=1 \let\foo=\relax \countdef\cc=7 \dimendef\dd=3
\let\bg={ \let\eg=} \chardef\ch=65 \toksdef\tk=4 \skipdef\sk=5
\def\mac{M}\let\mm=\mac
\tracingcommands=2
\foo
\cc=1
\dd=1pt
\setbox0=\hbox{\bg\eg}
\mm
\ch
\tk={x}
\sk=1pt
\let\pp=\par \pp
\tracingcommands=0
\end
